\chapter{Nelson--Siegel--Svensson modelling}\label{ch:nss}
\label{sec:nss-background}\label{sec:nss-equations}
Discrete maturity observations motivate smooth representation. The Nelson--Siegel model combines level, slope-related and curvature contributions; Svensson adds another curvature contribution, forming Nelson--Siegel--Svensson (NSS) \parencite{nelson1987,svensson1994}. Parameters specify a model member:
\(\theta=(\beta_0,\beta_1,\beta_2,\beta_3,\tau_1,\tau_2)\):
the \(\beta\) coefficients are decimal rates and positive decay scales \(\tau_1,\tau_2\) are years. A factor loading is a function of maturity multiplying a coefficient. For either decay scale \(\tau\), dimensionless loadings define the generic decimal rate \(r_{\mathrm{NSS}}\):
\begin{align}
L_1(T;\tau)&=\frac{1-e^{-T/\tau}}{T/\tau},&
L_2(T;\tau)&=L_1(T;\tau)-e^{-T/\tau},\label{eq:loadings}\\
r_{\mathrm{NSS}}(T;\theta)&=\beta_0+\beta_1L_1(T;\tau_1)
+\beta_2L_2(T;\tau_1)\notag\\
&\quad+\beta_3L_2(T;\tau_2).\label{eq:nss}
\end{align}
\label{sec:nss-limits}
Asymptotic limits, the values approached at boundaries, follow from the exponential expansions in \cref{sec:app-nss}:
\begin{equation}\label{eq:nss-limits}
r_{\mathrm{NSS}}(0^+)=\beta_0+\beta_1,\qquad
r_{\mathrm{NSS}}(\infty)=\beta_0.
\end{equation}
Thus \(\beta_0\) sets the long limit, \(\beta_1\) the short offset, and \(\beta_2,\beta_3\) overlapping hump/trough contributions; decay scales locate the individual loadings, not necessarily turning points of their sum. These are descriptive roles, not structurally unique economic factors.

\label{sec:nss-stability}
For dimensionless \(x=T/\tau\), subtracting nearly equal rounded numbers in \(1-e^{-x}\) causes cancellation. The implementation uses \texttt{expm1}, accurately evaluating \(e^{-x}-1\), and separate fifth-degree loading expansions below \(x=10^{-3}\) \parencite{higham2002}.

\section{Calibration and interpretation}\label{sec:calibration}
Calibration selects parameters fitting observations. For \(n\) observations \((u_\ell,z_\ell)\), indexed by \(\ell\), maturities are years and rates decimal. Residuals \(e_\ell\) are observed minus fitted:
\begin{equation}\label{eq:residuals}
e_\ell=z_\ell-r_{\mathrm{NSS}}(u_\ell;\theta),\qquad
S(\theta)=\sum_{\ell=1}^n e_\ell^2.
\end{equation}
Least squares minimises \(S\), measured in squared decimal rates. At the selected fit \(\widehat\theta\), root mean squared error is
\begin{equation}\label{eq:rmse}
\mathrm{RMSE}=\sqrt{S(\widehat\theta)/n},\qquad
\mathrm{RMSE}_{\bp}=10^4\mathrm{RMSE}.
\end{equation}
\label{sec:identification}
Nonlinear optimisation searches for parameters minimising \(S\); local minima minimise it within a neighbourhood, without ensuring global optimality. Six deterministic starts reduce starting-point dependence (\cref{sec:optimiser}). Excellent fit need not uniquely identify parameters \parencite{gilli2010}: a constant curve leaves decay scales undetermined. If \(\tau_1=\tau_2\), only the sum \(\beta_2+\beta_3\) enters their shared loading. Near equality can therefore produce unstable coefficients even when fitted rates change little.
\label{sec:nss-meaning}
The generic rate \(r_{\mathrm{NSS}}\) acquires zero-rate meaning only through the assumption in \cref{sec:zero-par}.
